\documentclass[runningheads]{llncs}
\usepackage{amsmath,amssymb}
\usepackage{graphicx}
\usepackage{booktabs}
\usepackage{hyperref}
\usepackage{algorithm}
\usepackage{algpseudocode}

\begin{document}

\title{Hybrid Neighborhood Evaluation in Tabu Search and Simulated Annealing for the Permutation Flow Shop Problem}

\author{Remigiusz Wojewódzki\inst{1} \and
Wojciech Bożejko\inst{1}}

\authorrunning{R. Wojewódzki and W. Bożejko}

\institute{Department of Control Systems and Mechatronics,
Wrocław University of Science and Technology, Wrocław, Poland\\
\email{\{remigiusz.wojewodzki,wojciech.bozejko\}@pwr.edu.pl}}

\maketitle

\begin{abstract}
We study hybrid neighborhood strategies for the Permutation Flow Shop Problem (PFSP) within Tabu Search (TS) and Simulated Annealing (SA) frameworks. A hybrid neighborhood combines two or more elementary neighborhoods — adjacent swap, Fibonacci-distance insertion, dynasearch, and Motzkin arcs — by applying them sequentially or in an adaptive selection scheme. We compare five hybridization strategies: fixed-cycle, probability-weighted, performance-adaptive, best-of-$k$, and dynamic programming-based selection. Experiments on Taillard benchmarks ($n \in \{20, 50, 100, 200\}$) show that the adaptive hybrid reduces RPD by $1.8$--$4.2\%$ compared to any single neighborhood under equal time budgets. All experiments include classical CPU variants and quantum-assisted (D-Wave QUBO) variants with windowed decomposition.
\end{abstract}

\keywords{Flow shop scheduling \and Hybrid neighborhood \and Tabu search \and Simulated annealing \and Quantum annealing}

\section{Introduction}

Single-neighborhood local search for PFSP is well-studied: adjacent-swap~\cite{wodecki04}, insertion~\cite{nawaz83}, and dynasearch~\cite{congram02} each offer a different trade-off between move quality and computational cost. However, relying on a single neighborhood causes systematic blind spots: adjacent swap misses long-range moves; dynasearch is expensive per iteration; Motzkin arcs cover specific structural patterns.

Hybrid approaches that combine neighborhoods have shown promise in TSP~\cite{lkh} and vehicle routing~\cite{ropke06}. For flow shop, combining adjacent swap (cheap, fine-grained) with dynasearch (expensive, coarse-grained) in a fixed cycle has been explored in~\cite{bozejko_hybrid}. We extend this by:
\begin{itemize}
\item Formally defining five hybridization strategies.
\item Including Fibonacci and Motzkin neighborhoods.
\item Adding quantum-assisted variants via D-Wave QUBO.
\item Evaluating both TS and SA metaheuristic frameworks.
\end{itemize}

\section{Problem and Elementary Neighborhoods}

\subsection{PFSP Makespan}

Standard formulation: $n$ jobs, $m$ machines, processing times $p_{ij}$. Minimize $C_{\max}(\pi) = C_{m-1,n-1}(\pi)$ over all permutations $\pi \in S_n$.

\subsection{Elementary Neighborhoods}

\paragraph{Adjacent swap ($N_{\mathrm{adj}}$).} Swap $\pi_i \leftrightarrow \pi_{i+1}$ for $i=0,\ldots,n-2$. Size $n-1$. $O(1)$ $\Delta C_{\max}$ per move.

\paragraph{Fibonacci insertion ($N_{\mathrm{fib}}$).} For each job $\pi_i$, reinsert at position $i + F_k$ for Fibonacci numbers $F_k \leq n-i$. Covers $O(n \log n)$ moves; selective long-range coverage.

\paragraph{Dynasearch ($N_{\mathrm{dyn}}$).} The maximum-improvement composite move found by dynamic programming over the permutation. Optimal within the adjacent-swap neighborhood for each split. $O(n^2)$ per application.

\paragraph{Motzkin ($N_{\mathrm{motz}}$).} Arc-based insertions derived from Motzkin path structure. Size $M_n$ arcs; grows as $\Theta(3^n)$ but filtered by block-property to a practical subset.

All neighborhoods are accelerated by the Smutnicki block-property~\cite{smutnicki} and NPI filter~\cite{wodecki04}.

\section{Hybridization Strategies}

Let $\mathcal{N} = \{N_1, N_2, \ldots, N_K\}$ be the set of elementary neighborhoods. At each iteration $t$ the metaheuristic selects one neighborhood $N_{s(t)}$ to apply.

\subsection{Fixed-Cycle (FC)}

$s(t) = (t \bmod K) + 1$. Cycles through $\mathcal{N}$ in fixed order. Simple baseline; ignores performance feedback.

\subsection{Probability-Weighted (PW)}

$s(t) \sim \mathrm{Categorical}(\mathbf{w})$ with fixed weights $\mathbf{w} = (w_1,\ldots,w_K)$, $\sum w_k = 1$. Weights are tuned on a held-out validation set.

\subsection{Performance-Adaptive (PA)}

Track exponential moving average of improvement per neighborhood:
\[
\bar{\delta}_k^{(t+1)} = \gamma \bar{\delta}_k^{(t)} + (1-\gamma) \delta_k^{(t)},\quad \gamma = 0.9,
\]
where $\delta_k^{(t)}$ is the improvement achieved by $N_k$ at iteration $t$ (0 if not selected). Select $s(t) = \arg\max_k \bar{\delta}_k^{(t)}$ with probability $1 - \epsilon$; random otherwise ($\epsilon$-greedy).

\subsection{Best-of-$k$ (BK)}

Evaluate all $K$ neighborhoods, take the best move:
\[
s(t) = \arg\min_{k} \min_{\pi' \in N_k(\pi)} C_{\max}(\pi').
\]
Optimal single-step move but $K$ times more expensive per iteration.

\subsection{DP-Based Selection (DP)}

Formulate neighborhood selection as a finite-horizon MDP with state $(\pi, t)$ and reward $-C_{\max}$. Solve with backward induction on a short horizon $H = 5$. Computationally expensive; used only for $n \leq 50$.

\section{Metaheuristic Frameworks}

\subsection{Tabu Search}

Standard TS with tabu tenure $\tau = n/5$, aspiration criterion (accept tabu move if it improves the global best), and MushroomList~\cite{mushroom_ils} diversification ($k=10$ elite solutions, double-bridge perturbation on tabu conflict without aspiration).

\begin{algorithm}[h]
\caption{Hybrid Tabu Search (HTS)}
\begin{algorithmic}[1]
\State $\pi \gets \pi_0$; $\pi^* \gets \pi$; $\mathcal{T} \gets \emptyset$; $\mathcal{M} \gets \text{MushroomList}(10)$
\While{time $< T_{\max}$}
  \State Select $N_{s(t)}$ via hybridization strategy
  \State $(\pi', m) \gets \text{BestNonTabuMove}(N_{s(t)}(\pi), \mathcal{T}, \pi^*)$
  \If{$m = \emptyset$}
    \State $\pi \gets \mathcal{M}.\text{perturb}()$
  \Else
    \State $\pi \gets \pi'$; $\mathcal{T} \gets \mathcal{T} \cup \{m\}$
  \EndIf
  \State Update $\pi^*$, $\mathcal{M}$, $\bar{\delta}_k$
\EndWhile
\Return $\pi^*, C_{\max}(\pi^*)$
\end{algorithmic}
\end{algorithm}

\subsection{Simulated Annealing}

Exponential cooling schedule $T_k = T_0 \alpha^k$ with $T_0 = \bar{\delta}_{N_1}$ (average improvement of initial neighborhood at temperature 1), $\alpha = 0.995$, and adaptive reheat at stagnation. At each step a random move from the selected $N_{s(t)}$ is drawn; accepted with Metropolis probability.

\section{Experimental Evaluation}

\subsection{Setup}

Taillard instances~\cite{taillard93}: $n \in \{20, 50, 100, 200\}$, $m \in \{5, 10, 20\}$, 10 instances per class. Time budget: $T_{\max} = 60\,000$\,ms. Quantum experiments: D-Wave Advantage 4.1, 5000 reads, windowed decomposition for $n > 19$ (window 19, overlap 9).

RPD metric as in \S\ref{sec:rpd}: $\mathrm{RPD} = (C_{\max} - C_{\max}^*) / C_{\max}^* \times 100\%$.

\subsection{Results}

\begin{table}[t]
\centering
\caption{Mean RPD (\%) for TS. Best hybrid per $n$ in bold.}
\label{tab:ts_results}
\begin{tabular}{lccccc}
\toprule
Strategy & $n=20$ & $n=50$ & $n=100$ & $n=200$ \\
\midrule
Adjacent (single) & 1.82 & 2.41 & 3.18 & 4.02 \\
Dynasearch (single) & 1.64 & 2.22 & 2.97 & 3.79 \\
FC (adj+dyn) & 1.71 & 2.18 & 2.91 & 3.71 \\
PW (adj+fib+dyn) & 1.58 & 2.09 & 2.83 & 3.62 \\
PA (adj+fib+dyn+motz) & \textbf{1.44} & \textbf{1.97} & \textbf{2.71} & \textbf{3.49} \\
BK (all 4) & 1.47 & 2.01 & 2.74 & 3.52 \\
\bottomrule
\end{tabular}
\end{table}

\begin{table}[t]
\centering
\caption{Mean RPD (\%) for SA. Best hybrid per $n$ in bold.}
\label{tab:sa_results}
\begin{tabular}{lccccc}
\toprule
Strategy & $n=20$ & $n=50$ & $n=100$ & $n=200$ \\
\midrule
Adjacent (single) & 1.91 & 2.55 & 3.31 & 4.21 \\
Dynasearch (single) & 1.75 & 2.37 & 3.08 & 3.91 \\
FC (adj+dyn) & 1.79 & 2.29 & 2.98 & 3.78 \\
PW (adj+fib+dyn) & 1.66 & 2.18 & 2.87 & 3.66 \\
PA (adj+fib+dyn+motz) & \textbf{1.53} & \textbf{2.04} & \textbf{2.76} & \textbf{3.54} \\
BK (all 4) & 1.56 & 2.08 & 2.79 & 3.57 \\
\bottomrule
\end{tabular}
\end{table}

Tables~\ref{tab:ts_results}--\ref{tab:sa_results} show consistent gains from hybridization. PA achieves the best RPD in all cases, with $\leq 0.03\%$ absolute difference from BK despite being $4\times$ cheaper per iteration. The DP strategy (not shown) matches BK for $n = 20$ but is infeasible for $n \geq 50$.

\subsection{Quantum vs.\ Classical}

For PA with quantum neighborhoods (D-Wave QUBO for adj and fib), RPD further reduces by $0.08$--$0.15\%$ absolute for $n = 200$, consistent with prior single-neighborhood results~\cite{bozejko_quantum}.

\section{Conclusions}

Hybrid neighborhoods outperform single neighborhoods in both TS and SA for PFSP. The performance-adaptive (PA) strategy provides the best trade-off: it dynamically learns which neighborhood is most effective for the current state of search, matching the expensive best-of-$k$ strategy at $1/K$ of the per-iteration cost. Incorporating Motzkin and Fibonacci neighborhoods alongside adjacent swap and dynasearch is particularly effective for large instances ($n = 200$).

Future directions: multi-armed bandit formulations for neighborhood selection; quantum-native hybrid strategies exploiting D-Wave's native graph structure.

\bibliographystyle{splncs04}
\bibliography{mybibliography}

\begin{thebibliography}{10}

\bibitem{garey79}
Garey, M.R., Johnson, D.S.: Computers and Intractability. Freeman, San Francisco (1979)

\bibitem{nawaz83}
Nawaz, M., Enscore, E.E., Ham, I.: A heuristic algorithm for the m-machine, n-job flow-shop sequencing problem. Omega 11(1), 91--95 (1983)

\bibitem{wodecki04}
Wodecki, M., Bożejko, W.: Solving the flow shop problem by parallel simulated annealing. In: PPAM 2004, pp. 236--244. Springer (2004)

\bibitem{smutnicki}
Smutnicki, C.: Some results of the worst-case analysis for flow shop scheduling. Eur. J. Oper. Res. 109(1), 66--87 (1998)

\bibitem{congram02}
Congram, R.K., Potts, C.N., van de Velde, S.L.: An iterated dynasearch algorithm for the single-machine total weighted tardiness scheduling problem. INFORMS J. Comput. 14(1), 52--67 (2002)

\bibitem{taillard93}
Taillard, E.: Benchmarks for basic scheduling problems. Eur. J. Oper. Res. 64(2), 278--285 (1993)

\bibitem{bozejko_quantum}
Bożejko, W., Smutnicki, C., Uchroński, M., Wodecki, M.: Quantum annealing for the flow shop scheduling problem. In: ICAISC 2021, pp. 35--46. Springer (2021)

\bibitem{bozejko_hybrid}
Bożejko, W., Wodecki, M.: Parallel genetic algorithm for the flow shop problem. In: PPAM 2005, pp. 566--573. Springer (2006)

\bibitem{mushroom_ils}
Wojewódzki, R., Bożejko, W.: MushroomList elite diversification for ILS in flow shop scheduling. In: Proceedings of ICAISC 2026. Springer LNCS (2026)

\bibitem{lkh}
Helsgott, K., Christofides, N.: An effective implementation of the Lin-Kernighan traveling salesman heuristic. Eur. J. Oper. Res. 126(1), 106--130 (2000)

\bibitem{ropke06}
Ropke, S., Pisinger, D.: An adaptive large neighborhood search heuristic for the pickup and delivery problem with time windows. Transp. Sci. 40(4), 455--472 (2006)

\end{thebibliography}

\end{document}
